\chapter{Comunicación entre servicios}
\label{cap:AnexoC}

Este anexo documenta los protocolos y contratos de comunicación entre los servicios del sistema: la API REST entre frontend y backend, el canal WebSocket para mensajes en tiempo real, y la llamada HTTP del backend al motor de IA.

\section{API REST: frontend \texorpdfstring{$\rightarrow$}{->} backend}

La base de la URL es configurable mediante la variable de entorno \texttt{NEXT\_PUBLIC\_API\_URL} (valor por defecto: \texttt{http://localhost:4000/api}). Las rutas protegidas requieren la cabecera \texttt{Authorization: Bearer <token>}, donde \texttt{<token>} es un JWT emitido por Phoenix al iniciar sesión con validez de 30 días.

\subsection{Autenticación}

\begin{table}[h]
\centering
\caption{Endpoints de autenticación.}
\label{tab:api-auth}
\begin{tabular}{llp{7cm}}
\hline
\textbf{Método} & \textbf{Ruta} & \textbf{Descripción} \\
\hline
POST & \texttt{/auth/register}       & Registra un nuevo usuario. Envía email de confirmación. \\
POST & \texttt{/auth/login}          & Autentica al usuario; devuelve \texttt{\{user, token\}}. \\
GET  & \texttt{/auth/confirm/:token} & Confirma el email del usuario a partir del enlace recibido. \\
PUT  & \texttt{/users/locale}        & Actualiza el idioma preferido del usuario. \emph{Protegido.} \\
\hline
\end{tabular}
\end{table}

El token se almacena en \texttt{localStorage} bajo la clave \texttt{cie10\_auth} como objeto \texttt{\{user, token\}}.

\subsection{Gestión de conversaciones}

\begin{table}[h]
\centering
\caption{Endpoints de conversaciones.}
\label{tab:api-conv}
\begin{tabular}{llp{7cm}}
\hline
\textbf{Método} & \textbf{Ruta} & \textbf{Descripción} \\
\hline
GET    & \texttt{/conversations?user\_id=<uuid>}                  & Lista conversaciones activas del usuario. \emph{Protegido.} \\
GET    & \texttt{/conversations/trash?user\_id=<uuid>}            & Lista conversaciones en papelera. \emph{Protegido.} \\
DELETE & \texttt{/conversations/:id}                             & Borrado lógico (soft delete). \emph{Protegido.} \\
PUT    & \texttt{/conversations/:id/restore}                     & Restaura conversación de la papelera. \emph{Protegido.} \\
\hline
\end{tabular}
\end{table}

\subsection{Catálogo CIE-10-ES}

\begin{table}[h]
\centering
\caption{Endpoints del catálogo CIE-10-ES (públicos, solo lectura).}
\label{tab:api-cie10}
\begin{tabular}{llp{7cm}}
\hline
\textbf{Método} & \textbf{Ruta} & \textbf{Descripción} \\
\hline
GET & \texttt{/cie10/search?q=<query>\&limit=<n>\&type=<tipo>} & Búsqueda en texto libre. \texttt{type} puede ser \texttt{diagnosis}, \texttt{procedure} o \texttt{chemical}. \\
GET & \texttt{/cie10/codes/:code}          & Devuelve un código concreto con su descripción y metadatos. \\
GET & \texttt{/cie10/codes/:code/children} & Devuelve los códigos hijos en la jerarquía CIE-10. \\
\hline
\end{tabular}
\end{table}

\subsection{Traducciones}

\begin{table}[h]
\centering
\caption{Endpoints de internacionalización (públicos).}
\label{tab:api-i18n}
\begin{tabular}{llp{7cm}}
\hline
\textbf{Método} & \textbf{Ruta} & \textbf{Descripción} \\
\hline
GET & \texttt{/translations/:locale} & Devuelve el mapa de cadenas de texto para \texttt{es} o \texttt{en}. \\
\hline
\end{tabular}
\end{table}

\section{Canal WebSocket: frontend \texorpdfstring{$\rightarrow$}{->} backend}

La conexión WebSocket se establece contra \texttt{ws://localhost:4000/socket} (configurable con \texttt{NEXT\_PUBLIC\_WS\_URL}) usando la librería \texttt{phoenix\_socket}. El token JWT se pasa como parámetro de conexión para autenticar la sesión.

Cada conversación se subscribe a un canal propio con el topic \texttt{conversation:<conversation\_id>}, donde \texttt{conversation\_id} es un UUID generado por el cliente.

\subsection{Mensajes entrantes (cliente \texorpdfstring{$\rightarrow$}{->} servidor)}

\begin{table}[h]
\centering
\caption{Mensajes enviados por el frontend al canal WebSocket.}
\label{tab:ws-in}
\begin{tabular}{lp{4.5cm}p{5.5cm}}
\hline
\textbf{Evento} & \textbf{Payload} & \textbf{Efecto} \\
\hline
\texttt{join}           & —                                                & Une al cliente al canal; devuelve el historial de la conversación. \\
\texttt{send\_message}  & \texttt{\{content: string\}}                     & Guarda el mensaje y lo difunde al canal. \\
\texttt{analyze\_report}& \texttt{\{report\_text: string\}}                & Persiste el informe y lanza la tarea asíncrona de análisis IA. \\
\texttt{validate\_code} & \texttt{\{code\_id, cie10\_code\}}               & Emite el comando \texttt{ValidateCode} al agregado. \\
\texttt{reject\_code}   & \texttt{\{code\_id, cie10\_code, reason\}}       & Emite el comando \texttt{RejectCode} al agregado. \\
\texttt{suggest\_code}  & \texttt{\{selected\_text, suggested\_code\}}     & Almacena una sugerencia manual del usuario. \\
\hline
\end{tabular}
\end{table}

\subsection{Mensajes salientes (servidor \texorpdfstring{$\rightarrow$}{->} cliente)}

\begin{table}[h]
\centering
\caption{Mensajes emitidos por el backend hacia el cliente WebSocket.}
\label{tab:ws-out}
\begin{tabular}{lp{4.5cm}p{5cm}}
\hline
\textbf{Evento} & \textbf{Payload principal} & \textbf{Cuándo se emite} \\
\hline
\texttt{joined}                  & \texttt{\{history\}}                               & Al unirse al canal con éxito. \\
\texttt{new\_message}            & \texttt{\{message\_id, content, user\_id, timestamp\}} & Cuando cualquier participante envía un mensaje. \\
\texttt{analysis\_started}       & \texttt{\{message\_id\}}                           & Inmediatamente tras recibir \texttt{analyze\_report}. \\
\texttt{analysis\_card\_received}& \texttt{\{card\_id, card\_type, content, message\_id\}} & Por cada una de las tres tarjetas generadas. \\
\texttt{analysis\_complete}      & \texttt{\{message\_id, predicted\_codes\}}         & Al completarse la predicción IA. \\
\texttt{code\_validated}         & \texttt{\{code\_id, cie10\_code\}}                 & Tras confirmar la validación en el agregado. \\
\texttt{code\_rejected}          & \texttt{\{code\_id, cie10\_code, reason\}}         & Tras confirmar el rechazo en el agregado. \\
\texttt{analysis\_failed}        & \texttt{\{message\_id, error\}}                    & Si el motor de IA devuelve un error. \\
\hline
\end{tabular}
\end{table}

Cada análisis genera exactamente tres tarjetas (\texttt{summary}, \texttt{codes}, \texttt{recommendations}) que se emiten como eventos \texttt{analysis\_card\_received} independientes antes del evento \texttt{analysis\_complete}.

\section{HTTP interno: backend \texorpdfstring{$\rightarrow$}{->} motor de IA}

El backend se comunica con el motor de IA mediante HTTP usando la librería Elixir \texttt{Req}. La URL base es \texttt{http://ai\_engine:8000} dentro de la red Docker, o \texttt{http://localhost:8000} en desarrollo sin Docker.

\subsection{Endpoint de predicción}

\begin{description}
    \item[Ruta] \texttt{POST /predict}
    \item[Cuerpo de la petición]~
\end{description}

\begin{verbatim}
{
  "text": "<informe clínico en español>",
  "engine": "bert" | "dict"   // opcional; por defecto "bert"
}
\end{verbatim}

\begin{description}
    \item[Respuesta] Objeto JSON con el campo \texttt{cards}, lista de tres objetos:
\end{description}

\begin{verbatim}
{
  "cards": [
    {
      "type": "summary",
      "content": "Informe analizado (349 palabras). 11 código(s) CIE-10 identificado(s)..."
    },
    {
      "type": "codes",
      "content": [
        {
          "code":        "E11",
          "description": "Diabetes mellitus tipo 2",
          "reason":      "Capítulo IV — confianza 95.3%",
          "confidence":  0.953,
          "engine":      "bert"
        }
      ]
    },
    {
      "type": "recommendations",
      "content": "Revisar y confirmar los códigos asignados con el equipo médico..."
    }
  ]
}
\end{verbatim}

La llamada se realiza dentro de una \texttt{Task.start} asíncrona, por lo que el backend no bloquea el canal WebSocket durante la inferencia. Cuando la respuesta llega, el backend emite el comando \texttt{ReceiveAIPrediction}, que genera los eventos correspondientes y activa los proyectores que persisten las tarjetas y los códigos predichos.

\subsection{Endpoint de estado}

\begin{description}
    \item[Ruta] \texttt{GET /}
    \item[Respuesta] \texttt{\{status: "ok", model\_loaded: bool, dict\_loaded: bool\}}
\end{description}

Permite al backend verificar que el motor de IA está disponible y que el modelo ha terminado de cargarse en memoria antes de aceptar peticiones de análisis.

\section{Mecanismo de resiliencia}

Si la llamada al motor de IA tiene éxito pero el despacho del comando \texttt{ReceiveAIPrediction} mediante Commanded falla, el backend dispone de un mecanismo de escritura directa en base de datos (\texttt{persist\_prediction\_direct/3}) que garantiza que las tarjetas y los códigos predichos se guardan igualmente. Esto asegura que el usuario recibe una respuesta aunque el \emph{event store} experimente un problema transitorio.
